%%%%%%%%%%%%%%%%%%%%%%%%%%%%%%%%%%%%%%%%%%%%%%%%%%%%%%%%%%%%%%%%%%%%%%%%
%
%  PREAMBLE.TEX
%  ------------
%
%	 This file contains settings and definitions of new commands.
%
%%%%%%%%%%%%%%%%%%%%%%%%%%%%%%%%%%%%%%%%%%%%%%%%%%%%%%%%%%%%%%%%%%%%%%
\documentclass[11pt,preprint]{aastex}
% \usepackage[dvips]{graphicx}
% \setlength{\parindent}{0.5cm}
% \setlength{\parskip}{0.2cm}
% \setlength{\textwidth}{6in}
% \setlength{\textheight}{8.5in}
% \setlength{\headsep}{0cm}
% \setlength{\headheight}{0cm}
% \setlength{\topskip}{0cm}
% \setlength{\topmargin}{0cm}
% \setlength{\evensidemargin}{0in}
% \setlength{\oddsidemargin}{0in}
% \setlength{\footskip}{0.5in}
% \setlength{\fboxsep}{0.5cm}
% \setlength{\fboxrule}{0.4mm}
%\addtolength{\textheight}{-\footskip}
\newcommand{\beq}[1]{\begin{equation} #1 \end{equation}}
\newcommand{\beqa}[1]{\begin{eqnarray} #1 \end{eqnarray}}
\newcommand{\greekbf}[1]{\mbox{\boldmath$#1$}}
\newcommand{\Item}{\hspace*{0.7cm}$\bullet$\hspace*{0.35cm}}
\newcommand{\Itemcirc}{\hspace*{0.7cm}$\circ$\hspace*{0.35cm}}
\newcommand{\smItem}{$\bullet$\hspace*{0.35cm}}
\newcommand{\It}{\hspace*{1.5cm}}
\newcommand{\T}{\hspace*{0.5cm}}
\newcommand{\makeline}[1]{\underline{\hspace*{#1}}}
\newcommand{\mathnorm}[1]{\mbox{\normalsize$#1$}}
\newcommand{\mathbig}[1]{\mbox{\large$#1$}}
\newcommand{\mathBig}[1]{\mbox{\Large$#1$}}
\newcommand{\mathsmall}[1]{\mbox{\small$#1$}}
\newcommand{\mathsmaller}[1]{\mbox{\footnotesize$#1$}}
\newcommand{\mathsmallest}[1]{\mbox{\scriptsize$#1$}}
\newcommand{\mathtiny}[1]{\mbox{\tiny$#1$}}
\newcommand{\mathbox}[1]{\fbox{$\displaystyle #1 $}}
\newcommand{\switchfonts}{\usefont{OT1}{cmr}{m}{it}}
\newcommand{\lie}{\switchfonts \mbox{\symbol{36}} \normalfont}
\newcommand{\deriv}[2]{\frac{ d #1 }{ d #2 }}
\newcommand{\pderiv}[2]{\frac{ \partial #1 }{ \partial #2 }}
\newcommand{\bfbar}[1]{\mathbf{\bar{#1}}}
\newcommand{\qbar}{\bfbar{q}}
\newcommand{\ibanez}{J. $\mathrm{M}^{\underline{a}}$ Ib\'{a}\~{n}ez}
\newcommand{\marti}{J. $\mathrm{M}^{\underline{a}}$ Mart\'{i}}
\newcommand{\bftt}[1]{\textbf{\texttt{#1}}}
\newcommand{\maxwell}[2]{{\,}^*\!F^{#1 #2}}

\title{Derivation of the Divergence Cleaning Equations of Motion}
\author{Scott Noble, Josh Faber, Bruno Mundim}
\date{}
\begin{document}

I'll use \cite{anton}  and \cite{penner}.

We start with the modified form of Maxwell's equations in covariant form with the divergence cleaning field, $\psi$:
\beq{
\nabla_\mu \left( \maxwell{\mu}{\nu} + g^{\mu \nu} \psi \right) =  \kappa n^\nu \psi
\label{general-div-cleaning-eq}
}
which comes from \cite{penner} except we correct the sign of the RHS.  Also, note that when $\kappa > 0$, $\nabla^\mu \nabla_\mu \psi = \kappa \nabla_\mu n^\mu \psi$ is 
a damped wave equation.   We will thus use $\kappa > 0$ as $\partial_t \psi$ is proportional to the divergence of the magnetic field which we wish to 
drive to zero. 
We can simplify the original part of Maxwell's equations:
\beqa{
\nabla_\mu   \maxwell{\mu}{\nu}  & = &\partial_\mu \maxwell{\mu}{\nu}+{\Gamma^\mu}_{\lambda\mu} 
\maxwell{\lambda}{\nu}+{\Gamma^\nu}_{\lambda\mu}\maxwell{\mu}{\lambda}\\
& = &\frac{1}{\sqrt{-g}} \partial_\mu \left( \sqrt{-g}  \maxwell{\mu}{\nu} \right)  
+ {\Gamma^\nu}_{\mu \lambda} \maxwell{\mu}{\lambda}   \\
& = & \frac{1}{\sqrt{-g}} \partial_\mu \left( \sqrt{-g}  \maxwell{\mu}{\nu} \right)  
}
where the connection term vanished because $\maxwell{\mu}{\lambda}$ is anti-symmetric while ${\Gamma^\nu}_{\mu \lambda}$
is symmetric under permutation of its lower indices. 

Using equation~(18) from \cite{anton}, 
\beq{
\maxwell{\mu}{\nu} = \frac{1}{W} \left( u^\mu B^\nu - u^\nu B^\mu \right)  \label{maxwell-tensor}
}
where $B^\mu$ is a purely spatial vector and is the magnetic field w.r.t. the hypersurfaces normal observer, i.e. the one we want. 
Therefore, Maxwell's equations become (remembering that $B^t=0$ and $\maxwell{t}{t} = 0$):
\beqa{
\nabla_\mu  \maxwell{\mu}{j} & = & \frac{1}{\sqrt{-g}} \partial_\mu \left( \sqrt{-g}  \maxwell{\mu}{j}  \right)  \\
& = & \frac{1}{\sqrt{-g}} \left[ \partial_t \sqrt{-g} \frac{1}{W} \left( u^t B^j \right) + 
\partial_i \sqrt{-g} \frac{1}{W} \left( u^i B^j - u^j B^i \right) \right] \\
& = & \frac{1}{\alpha \sqrt{\gamma}} \left[ \partial_t \alpha \sqrt{\gamma} \frac{1}{\alpha u^t} \left( u^t B^j \right) + 
\partial_i \alpha \sqrt{\gamma} \frac{1}{\alpha u^t} \left( u^i B^j - u^j B^i \right) \right] \\
& = & \frac{1}{\alpha \sqrt{\gamma}} \left\{ \partial_t \sqrt{\gamma}  B^j + 
\partial_i  \sqrt{\gamma} \left[ \left(\alpha v^i - \beta^i\right) B^j - \left(\alpha v^j - \beta^j\right) B^i \right] \right\}
}
which is equation~(20) from \cite{anton} divided by $\alpha$. 
Note we have used the following identity in arriving at the last expression:
\beq{
\frac{u^i}{u^t} = \alpha v^i - \beta^i 
}

The divergence constraint comes from the time component of the Maxwell's equation:
\beqa{
0=\nabla_\mu  \maxwell{\mu}{t}  & = & \frac{1}{\sqrt{-g}} \partial_i \left( \sqrt{-g}  \maxwell{i}{t} \right)  \\
& = & \frac{1}{\sqrt{-g}}  \partial_i \sqrt{-g} \frac{1}{W} \left( u^i B^t - u^t B^i \right)  \\
& = & - \frac{1}{\alpha \sqrt{\gamma}} \partial_i \alpha \sqrt{\gamma} \frac{1}{\alpha u^t} u^t B^i  \\
& = & - \frac{1}{\alpha \sqrt{\gamma}} \partial_i \sqrt{\gamma} B^i  
}

In 3+1 form, the metric's inverse is defined as 
\beq{
g^{\mu \nu} = \left[ 
\begin{array}{cc} -\frac{1}{\alpha^2} &\frac{\beta^i}{\alpha^2} \\[0.25cm]
\frac{\beta^j}{\alpha^2}  & \gamma^{i j} - \frac{\beta^i \beta^j}{\alpha^2}
\end{array} \right] \quad .
\label{inverse-metric}
}
and 
\beq{
n_\mu = \left[-\alpha,0,0,0\right] \quad , \quad 
n^\mu = \frac{1}{\alpha} \left[1,-\beta^j\right]^T \quad , \label{normal-observer}
}

%   The $t$-component of the divergence cleaning term becomes 
%   \beqa{ 
%   - \frac{\kappa}{\alpha} \psi  = \nabla_\mu  g^{\mu t} \psi  & = & \frac{1}{\sqrt{-g}} \partial_\mu \left( \sqrt{-g} g^{\mu t} \psi \right)
%   + {\Gamma^t}_{\mu \lambda} g^{\lambda \mu} \psi \\
%    & = & - \frac{1}{\alpha \sqrt{\gamma}} \partial_t \left( \sqrt{\gamma}  \frac{\psi}{\alpha}  \right)
%   + \frac{1}{\alpha \sqrt{\gamma}} \partial_i \left( \sqrt{\gamma} \frac{\psi}{\alpha} \beta^i \right)
%   + {\Gamma^t}_{\mu \lambda} g^{\lambda \mu} \psi 
%   }
%   
%   Josh's connection identity is:
%   \beq{
%   \sqrt{-g} g^{\mu \nu} {\Gamma^\kappa}_{\mu \nu} = - \partial_\mu \left( \sqrt{-g} g^{\kappa \mu}  \right)
%   }
%   
%   The $j$-component of the divergence cleaning term becomes 
%   \beqa{ 
%    \frac{\kappa}{\alpha} \beta^j \psi  = \nabla_\mu  g^{\mu j} \psi  & = & \frac{1}{\sqrt{-g}} \partial_\mu \left( \sqrt{-g} g^{\mu j} \psi \right)
%   + {\Gamma^j}_{\mu \lambda} g^{\lambda \mu} \psi \\
%   & = & \frac{\psi}{\sqrt{-g}} \partial_\mu \left( \sqrt{-g} g^{\mu j} \right) + g^{\mu j}  \partial_\mu  \psi 
%   + {\Gamma^j}_{\mu \lambda} g^{\lambda \mu} \psi \\
%   & = & g^{\mu j}  \partial_\mu  \psi - \frac{\psi}{\sqrt{-g}} \sqrt{-g} g^{\mu \nu} {\Gamma^j}_{\mu \nu}  
%   + {\Gamma^j}_{\mu \lambda} g^{\lambda \mu} \psi \\
%   & = & g^{\mu j}  \partial_\mu  \psi - \psi g^{\mu \nu} {\Gamma^j}_{\mu \nu}  
%   + \psi g^{\lambda \mu} {\Gamma^j}_{\mu \lambda}  \\
%   & = & g^{\mu j}  \partial_\mu  \psi \label{psi-j-eq} \\
%   % & = &  \frac{1}{\alpha \sqrt{\gamma}} \partial_t \left( \sqrt{\gamma}  \frac{\psi}{\alpha} \beta^j  \right)
%   %+ \frac{1}{\alpha \sqrt{\gamma}} \partial_i \left[ \sqrt{\gamma} \frac{\psi}{\alpha} \left( \alpha^2 \gamma^{i j} - \beta^i \beta^j \right) \right]
%   %+ {\Gamma^j}_{\mu \lambda} g^{\lambda \mu} \psi 
%   }

% The connection terms look like:
% \beqa{
% {\Gamma^\nu}_{\mu \lambda} g^{\lambda \mu} & = &
% \frac{1}{2} g^{\lambda \mu} g^{\nu \kappa} \left( \partial_\mu g_{\lambda \kappa} + \partial_\lambda g_{\mu \kappa} 
% - \partial_\kappa g_{\mu \lambda} \right)\\
% & = & \frac{1}{2} g^{\lambda \mu} g^{\nu \kappa} \left( 2 \partial_\mu g_{\lambda \kappa} 
% - \partial_\kappa g_{\mu \lambda} \right)\\
% }

To simplify/expand the divergence cleaning term in Eq.~\ref{general-div-cleaning-eq}, we may take out the metric from its covariant derivative:
\beq{
\nabla_\mu g^{\mu \nu} \psi = g^{\mu \nu} \nabla_\mu \psi = g^{\mu \nu} \partial_\mu \psi 
}

The $t$-component is 
\beq{
\nabla_\mu g^{\mu t} \psi = g^{\mu t} \partial_\mu \psi = \frac{1}{\alpha^2} \left[ -\partial_t \psi + \beta^i \partial_i \psi \right]
}
which leads to the ultimate equation:
\beq{
\frac{1}{\alpha^2} \left[ \partial_t \psi - \beta^i \partial_i \psi \right] + \frac{1}{\alpha \sqrt{\gamma}} \partial_i \sqrt{\gamma} B^i = -\frac{\kappa}{\alpha} \psi  
}
\beq{
\rightarrow \partial_t \psi - \beta^i \partial_i \psi  + \frac{\alpha}{\sqrt{\gamma}} \partial_i \sqrt{\gamma} B^i = -\kappa \alpha \psi  
}
\beq{
\fbox{$ \displaystyle
\partial_t \psi + \partial_i \left( \alpha B^i - \psi \beta^i \right)  
=  \psi \left( -\kappa \alpha - \partial_i \beta^i  \right)
+ \sqrt{\gamma} B^i \partial_i \left( \frac{\alpha}{\sqrt{\gamma}} \right)\label{dphi-dt-eq}
$}
}

The $j$-component is 
\beq{
\nabla_\mu g^{\mu j} \psi =  g^{\mu j} \partial_\mu \psi 
= \frac{1}{\alpha^2} \left[ \beta^j \partial_t \psi + \left(\alpha^2 \gamma^{i j} - \beta^i \beta^j \right) \partial_i \psi \right]
=  -\frac{\kappa}{\alpha} \beta^j \psi   
}
\beq{
\rightarrow  \beta^j \partial_t \psi + \left(\alpha^2 \gamma^{i j} - \beta^i \beta^j \right) \partial_i \psi
=  -\kappa \alpha  \beta^j \psi  
}

The complete $j^\mathrm{th}$-component of the modified Maxwell's equation becomes
\beq{
\frac{1}{\alpha \sqrt{\gamma}} \left[ \partial_t \sqrt{\gamma}  B^j + 
\partial_i  \sqrt{\gamma} \left( u^i B^j - u^j B^i \right) \right]
+ g^{j \mu} \partial_\mu \psi 
 =  -\kappa \frac{\beta^j}{\alpha} \psi  
}
Inserting equation~\ref{dphi-dt-eq} :
\beqa{
&& \frac{1}{\alpha \sqrt{\gamma}} \left[ \partial_t \sqrt{\gamma}  B^j + 
\partial_i  \sqrt{\gamma} \left( u^i B^j - u^j B^i \right) \right]
- \frac{\beta^j}{\alpha^2}\partial_i\left( \alpha B^i - \psi\beta^i \right) 
+ g^{i j} \partial_i \psi \\
& = & -\kappa \frac{\beta^j}{\alpha} \psi  + \frac{\beta^j}{\alpha^2} \left\{ \kappa \alpha \psi  +\psi\partial_i\beta^i-\sqrt{\gamma}B^i\partial_i\left(\frac{\alpha}{\sqrt{\gamma}}\right) \right\} 
}
Multiplying through by $\sqrt{-g}=\alpha\sqrt{\gamma}$:
\beqa{
&&\partial_t \sqrt{\gamma}  B^j + 
\partial_i  \sqrt{\gamma} \left( u^i B^j - u^j B^i \right) 
- \frac{\sqrt{\gamma}\beta^j}{\alpha}\partial_i \left(\alpha B^i - \psi\beta^i \right) 
+\alpha\sqrt{\gamma} g^{i j} \partial_i \psi \\
& = & \frac{\sqrt{\gamma}\beta^j}{\alpha} \left\{\psi\partial_i\beta^i-\sqrt{\gamma}B^i\partial_i\left(\frac{\alpha}{\sqrt{\gamma}}\right)  \right\} 
}
Grouping spatial derivatives:
\beqa{
&&\partial_t \sqrt{\gamma}  B^j + 
\partial_i  \sqrt{\gamma} \left( u^i B^j - u^j B^i \right) 
+\frac{\sqrt{\gamma}\beta^j}{\alpha}\left[  \sqrt{\gamma}B^i\partial_i\left(\frac{\alpha}{\sqrt{\gamma}}\right)-\partial_i(\alpha B^i)\right]
+\nonumber\\
&&~~\alpha\sqrt{\gamma}\left( g^{i j} \partial_i \psi+\frac{\beta^j}{\alpha^2}[\partial_i(\psi\beta^i)-\psi\partial_i\beta^i]\right)
=  0
}
Simplifying the terms:
\beqa{
&&\partial_t \sqrt{\gamma}  B^j + 
\partial_i  \sqrt{\gamma} \left( u^i B^j - u^j B^i \right) 
-\frac{\sqrt{\gamma}\beta^j}{\alpha}\left[ \frac{\alpha}{\sqrt{\gamma}}\partial_i(\sqrt{\gamma}B^i)\right]
+\alpha\sqrt{\gamma}\left( g^{i j} +\frac{\beta^i\beta^j}{\alpha^2}\right)\partial_i\psi\nonumber
\\&&~~= 0
}
which reduces to:
\beqa{
&&\partial_t \sqrt{\gamma}  B^j + 
\partial_i  \sqrt{\gamma} \left( u^i B^j - u^j B^i \right) 
-\beta^j\partial_i(\sqrt{\gamma}B^i)
+\alpha\sqrt{\gamma} \gamma^{i j} \partial_i\psi=  0
}
or in conservative form:
\beq{
\fbox{$ \begin{array}{ll} 
 \partial_t \sqrt{\gamma}  B^j & + \partial_i  \sqrt{\gamma} \left[ u^i B^j - u^j B^i + \alpha \gamma^{i j} \psi 
- \beta^j B^i \right] 
\\[0.25cm] 
& = - \sqrt{\gamma} \left(  B^i \partial_i \beta^j  \right)
+ \psi \partial_i \left( \alpha \sqrt{\gamma} \gamma^{i j}  \right) 
\end{array} $}
}
%%%%%%%%%%%%%%%%%%%%%%%%%%%%%%%%%%%%%%%%%%%%%%%%%%%%%%%%%%%%%%%%%%%%%%
\begin{thebibliography}{9}

\bibitem[Penner(2011)]{penner} J.~Penner, \emph{PhD Thesis}, University of British Columbia (2011); 
\url{http://bh0.physics.ubc.ca/People/ajpenner/thesis/phd.pdf}. 

\bibitem[Ant\'{o}n et al.(2006)]{anton} L.~Ant\'{o}n, et al., \emph{ApJ}, \textbf{637}, 296-312 (2006).

\end{thebibliography}

\end{document}
























